% Uncooked resume template: the layout every converted resume uses.
% It passes check_layout.py (rules R1-R7, R10) and compiles on Overleaf as it is.
%
% Filling it from a PDF or Word resume:
% - Copy the person's words exactly. Only the layout changes: no rewording, no new words,
%   nothing dropped. Keep their section order and their own heading words.
% - Keep the preamble (everything before \begin{document}) unchanged.
% - Replace each {{placeholder}} (braces included) and each example address.
% - Repeat a block once per job, project or degree, with \vspace{6pt} between blocks.
%   Delete a line or section the person doesn't have; never fill it with a guess.
% - Escape LaTeX's special characters in their text:
%     & -> \&    % -> \%    $ -> \$    # -> \#    _ -> \_    { -> \{    } -> \}
%     ~ -> \textasciitilde{}    ^ -> \textasciicircum{}    \ -> \textbackslash{}
%   A line that starts with [ right after \\ or \item needs {} in front: \item {}[2024] ...
% - Date ranges use -- ("Mar 2024 -- Present"). Separators between fields are
%   \quad$\cdot$\quad on contact lines and $\cdot$ inside entries. Bullet characters
%   from the original go: itemize draws the bullets.
% - Links: \href{https://full-address}{address-without-https}. When the original showed only
%   a word ("LinkedIn") over a hidden address, show the address so parsers can read it.
\documentclass[10pt, letterpaper]{article}
\usepackage[margin=0.5in]{geometry}
\usepackage{iftex}
\ifpdftex
  \input{glyphtounicode}
  \pdfgentounicode=1
\fi
\usepackage[none]{hyphenat}
\usepackage{hyperref}
\usepackage{enumitem}
\usepackage{titlesec}
\usepackage{xcolor}

\hypersetup{
    colorlinks=true,
    linkcolor=black,
    urlcolor=blue,
}

\titleformat{\section}{\large\bfseries\uppercase}{}{0em}{}[\vspace{2pt}\titlerule]
\titlespacing{\section}{0pt}{8pt}{4pt}

\setlist[itemize]{leftmargin=*, noitemsep, topsep=0pt, parsep=0pt, itemsep=0.5pt}

\begin{document}
\pagestyle{empty}

% Contact: one line per line of the original, fields joined by \quad$\cdot$\quad.
\begin{center}
    {\LARGE \textbf{{{FULL NAME}}}} \\ \vspace{2pt}
    {{City, Country}} \quad$\cdot$\quad {{Phone}} \quad$\cdot$\quad \href{mailto:name@example.com}{name@example.com} \\
    \href{https://github.com/your-name}{github.com/your-name} \quad$\cdot$\quad \href{https://www.linkedin.com/in/your-name}{linkedin.com/in/your-name} \\
    {{Work authorization or other header line, if the original has one}}
\end{center}

\section*{Summary}
{{Summary paragraph, word for word.}}

% Skills: the person's own group labels. With no groups on the original, one plain list, no label.
\section*{Technical Skills}
\textbf{{{Group}}:} {{Skill}}, {{Skill}}, {{Skill}} \\
\textbf{{{Group}}:} {{Skill}}, {{Skill}}, {{Skill}}

% Experience: company line, then title line. Several roles at one company: one title line
% (with its own dates and bullets) per role under a single company line.
\section*{Experience}

\textbf{{{Company}}} \hfill {{Location}} \\
\textit{{{Job Title}}} \hfill {{Mon YYYY}} -- {{Mon YYYY or Present}}
\begin{itemize}
    \item {{Bullet, word for word.}}
    \item {{Bullet, word for word.}}
\end{itemize}

\vspace{6pt}
\textbf{{{Company}}} \hfill {{Location}} \\
\textit{{{Job Title}}} \hfill {{Mon YYYY}} -- {{Mon YYYY}}
\begin{itemize}
    \item {{Bullet, word for word.}}
\end{itemize}

\section*{Projects}

\textbf{{{Project or role}}} $\cdot$ \textit{{{Kind, e.g. Side Project}}} ({{Tools, if listed}}) \hfill {{Dates}}
\begin{itemize}
    \item {{Bullet, word for word.}}
\end{itemize}

% Education: stacked, one fact per line, no \hfill (rule R10). \vspace{6pt} between schools.
\section*{Education}
\textbf{{{School}}} \\
{{Degree}} \\
{{Location}} \quad$\cdot$\quad {{Graduation date}} \\
{{GPA, honours or other line, if the original has one}}

% Any other section (Certifications, Awards, Volunteering, Languages...): the original's
% heading, then its lines as entries like above, or as a list.
\section*{Certifications}
\begin{itemize}
    \item {{Certification}} $\cdot$ {{Issuer}} $\cdot$ {{Date}}
\end{itemize}

\end{document}
